The primary prediction received partial support. Discarding one nuisance direction consistently reduced orientation dependence, while the digit-accuracy gain was small and varied by seed. The single split, five model seeds, known template and synthetic fixed-phase stripes restrict the scope. All final comparisons used the same budget. At strong amplitude none of the five standard or projected runs reached the original assignment-change threshold at the last refresh. Convergence, other forms of background variation, and what the representation retains before the centroids are assigned remain open questions.

Two training confusion matrices admitted equally good cluster-name maps with different test ACC, including one strong-stripe projected model (62.14--62.64\%). We retained the fixed training-time matching convention. Flip rate and NMI use the original cluster IDs and therefore do not depend on this tie. A follow-up could apply the projection during pretraining and then during fine-tuning in a crossed design to test when the nuisance becomes established.
